\documentclass[11pt,a4paper]{article}
\usepackage[T1]{fontenc}
\usepackage{lmodern,amsmath,amssymb,amsthm,mathtools,microtype}
\usepackage[margin=27mm]{geometry}
\usepackage{booktabs,xurl,hyperref,fancyhdr}
\hypersetup{hidelinks,pdftitle={FR-12: A matching-indexed doubling counterexample},pdfauthor={George Stepaniants}}
\pagestyle{fancy}\fancyhf{}\fancyhead[L]{\small FR-12 / Hadamard enumeration}\fancyfoot[C]{\thepage}
\setlength{\headheight}{14pt}\setlength{\parskip}{4pt}\setlength{\emergencystretch}{2em}
\newtheorem{theorem}{Theorem}\newtheorem{lemma}{Lemma}\newtheorem{corollary}{Corollary}
\newcommand{\HH}{\mathcal H}
\title{\textbf{FR-12: A matching-indexed\\doubling counterexample}}
\author{George Stepaniants\\[3pt]\small Department of Computing and Mathematical Sciences\\\small California Institute of Technology\\\small Pasadena, California, USA}
\date{12 September 2026}
\begin{document}\maketitle
\begin{abstract}
Let $H(n)$ count labeled real Hadamard matrices of order $n$. An injective
construction gives $H(2m)\ge(2m-1)!!H(m)^2$. Iteration yields
$H(2^k)\ge 2^{2^k(k-1)(k-2)/8}$ for $k\ge2$.
Consequently no absolute constant $C$ satisfies $H(n)\le2^{Cn\log_2 n}$
at every order divisible by four. This refutes the exact counting target of
FR-12; it does not address existence at every admissible order.
\end{abstract}
\noindent\textbf{Assistance and independent review.}
Prepared with substantial assistance from Codex. A separate Codex agent
independently audited the complete supplied argument and reported PASS on
12 September 2026. The
\href{../../references/stepaniants-fr12-2026-09-12/REVIEW.md}{signed independent review}
identifies the exact archived source and scope. This is an informal AI-agent
audit, not external human peer review or formal verification. The
supplementary finite checker is not a universal proof certificate, and no
historical priority certification is claimed.

\section{Exact target}
Write
\[
 \HH_n=\{A\in\{-1,1\}^{n\times n}:AA^{\mathsf T}=nI_n\},
 \qquad H(n)=|\HH_n|.
\]
Rows and columns are labeled; signed permutations are \emph{not} quotiented
out. The target in FR-12~\cite{catalog}, attributed to
Ferber--Jain--Zhao~\cite{FJZ}, is whether an absolute $C>0$ gives
\begin{equation}\label{eq:target}
 H(n)\le 2^{Cn\log_2 n}\qquad(4\mid n).
\end{equation}
We count actual matrices, so no assumption about automorphism groups or
sizes of equivalence classes enters the argument.

\begin{lemma}[Injective doubling]\label{lem:doubling}
For every positive integer $m$,
\begin{equation}\label{eq:recurrence}
 H(2m)\ge \frac{(2m)!}{2^m m!}\,H(m)^2=(2m-1)!!\,H(m)^2.
\end{equation}
\end{lemma}
\begin{proof}
Let $\mathfrak M_{2m}$ denote the perfect matchings of the labeled positions
$\{1,\ldots,2m\}$. Given a matching, write each pair as $(u_i,v_i)$ with
$u_i<v_i$, and order its pairs by $u_1<\cdots<u_m$.
For $A,B\in\HH_m$, define a $2m\times2m$ sign matrix $C$ by
\begin{equation}\label{eq:construction}
 C_{u_i,\cdot}=(A_{i,\cdot},B_{i,\cdot}),\qquad
 C_{v_i,\cdot}=(A_{i,\cdot},-B_{i,\cdot}),\quad 1\le i\le m.
\end{equation}
Each row has squared norm $2m$. The inner product of the two rows in one
pair is $m-m=0$. For two different pairs $i\ne j$, all four possible inner
products vanish, since they have the form
\[
 \langle A_{i,\cdot},A_{j,\cdot}\rangle
 \mathbin{\pm}\langle B_{i,\cdot},B_{j,\cdot}\rangle=0.
\]
Thus $C\in\HH_{2m}$.

The map $(A,B,\mathcal M)\mapsto C$ is injective. Indeed, distinct rows of
$A$ are distinct vectors: equality of two rows would contradict their
orthogonality. It follows that, among the rows of $C$, equality of the
\emph{first $m$ coordinates} partitions the $2m$ row positions into exactly
the pairs of $\mathcal M$. Hence $C$ recovers the matching. Its canonical
pair ordering then recovers $A$ from these first halves. For each pair,
the last $m$ coordinates of the smaller-indexed row recover the corresponding
row of $B$. This also fixes every sign, so no residual multiplicity remains.

Finally, $|\mathfrak M_{2m}|=(2m)!/(2^m m!)$: list all positions in an order,
pair successive positions, and divide by the $2^m$ orders within pairs and
the $m!$ orders of the pairs. Counting the domain of the injection proves
\eqref{eq:recurrence}.
\end{proof}

\begin{theorem}[A super-$n\log n$ exponent]\label{thm:main}
For every integer $k\ge2$,
\begin{equation}\label{eq:bound}
 \boxed{H(2^k)\ge 2^{\,2^k(k-1)(k-2)/8}.}
\end{equation}
In particular, the assertion \eqref{eq:target} is false.
\end{theorem}
\begin{proof}
Hadamard matrices exist at every power of two, by applying
\eqref{eq:construction} recursively starting with an order-one sign matrix.
Consequently the logarithms below are defined. Put
\[
 a_k=2^{-k}\log_2 H(2^k).
\]
For $k\ge2$, let $m=2^{k-1}$, so $m$ is even. The matching factor satisfies
\[
 (2m-1)!!=\prod_{j=1}^m(2j-1)\ge m!
 \ge (m/2)^{m/2}.
\]
The last inequality retains just the largest $m/2$ factors in $m!$ and
bounds each from below by $m/2$. Taking logarithms of
\eqref{eq:recurrence} and dividing by $2m$ gives
\[
 a_k\ge a_{k-1}+\frac{\log_2m-1}{4}
       =a_{k-1}+\frac{k-2}{4}.
\]
Since $a_1\ge0$, summation yields
\[
 a_k\ge\frac14\sum_{j=2}^k(j-2)
      =\frac{(k-1)(k-2)}8,
\]
which proves \eqref{eq:bound}.
If \eqref{eq:target} held, applying it at $n=2^k$ would give $a_k\le Ck$.
This contradicts the displayed quadratic lower bound for sufficiently large
$k$. All these $n$ are divisible by four.
\end{proof}

\section{Checks, attribution, and scope}
The supplementary script checks the equivalent row-permutation formulation:
form $\left(\begin{smallmatrix}A&B\\A&-B\end{smallmatrix}\right)$, then
permute its $2m$ rows. Every output has exactly $2^m m!$ preimages in that
formulation. The matching-indexed construction in Lemma~\ref{lem:doubling}
removes this multiplicity explicitly.
\begin{center}
\begin{tabular}{rrrrr}\toprule
$m$ & $H(m)$ & Input triples & Distinct outputs & Fiber size\\\midrule
1 & 2 & 8 & 4 & 2\\
2 & 8 & 1536 & 192 & 8\\\bottomrule
\end{tabular}
\end{center}
Those finite cases, and exact integer checks of the iterated lower bound
through $k=10$, are not used to infer an asymptotic statement. The analytic
injection and its iteration establish the result for every $k$.

The construction uses only power-of-two orders. It neither resolves the
Hadamard existence conjecture nor supplies a matching upper bound on $H(n)$.
The original counting question and the published upper bounds remain
attributed to their sources. The accompanying audit states the observed
public branch scope and the limits of the duplicate/literature search.

\begin{thebibliography}{9}\raggedright
\bibitem{FJZ}
A. Ferber, V. Jain, and Y. Zhao,
\emph{On the number of Hadamard matrices via anti-concentration},
Combinatorics, Probability and Computing 31 (2022), 455--477.
\href{https://doi.org/10.1017/S0963548321000377}{doi:10.1017/S0963548321000377}.
Published Conjecture 1.3; Conjecture 1.4 in
\href{https://arxiv.org/abs/1808.07222}{arXiv:1808.07222v1}.
\bibitem{catalog}
\emph{Open Problems in Numerical Linear Algebra}, entry FR-12,
\texttt{ajt60gaibb/OpenProblemsInNLA}, snapshot
\texttt{1f22006bdaa4659fcaa0bb775a887685cd3cc566}.
\url{https://github.com/ajt60gaibb/OpenProblemsInNLA/blob/1f22006bdaa4659fcaa0bb775a887685cd3cc566/frames-and-matrix-designs/FR-12/README.md}.
\end{thebibliography}
\end{document}
